\documentclass{article}
\usepackage[margin=28mm]{geometry}
\usepackage{amsmath,amssymb}
\usepackage{xurl}
\usepackage[hidelinks]{hyperref}
% Shared commands for notes in docs/paper-gaps/.

\DeclareUnicodeCharacter{2080}{\ifmmode{}_{0}\else\textsubscript{0}\fi}
\DeclareUnicodeCharacter{2081}{\ifmmode{}_{1}\else\textsubscript{1}\fi}
\DeclareUnicodeCharacter{2082}{\ifmmode{}_{2}\else\textsubscript{2}\fi}
\DeclareUnicodeCharacter{2099}{\ifmmode{}_{n}\else\textsubscript{n}\fi}

\newcommand{\Lean}{\textsc{Lean}}
\ExplSyntaxOn
\NewDocumentCommand{\leanid}{m}{
  \ifmmode
    \textsf{\detokenize{#1}}
  \else
    \group_begin:
    \urlstyle{sf}
    \tl_set:Nn \l_tmpa_tl {#1}
    \tl_replace_all:Nnn \l_tmpa_tl {\_} {_}
    \exp_args:NV \path \l_tmpa_tl
    \group_end:
  \fi
}
\ExplSyntaxOff
\providecommand{\tr}{\operatorname{tr}}
\newcommand{\tnleanIssueBase}{https://github.com/LionSR/TNLean/issues/}
\newcommand{\tnleanPrBase}{https://github.com/LionSR/TNLean/pull/}
\newcommand{\tnleanDocBase}{https://sirui-lu.com/TNLean/docs/}
\newcommand{\ghissue}[1]{\href{\tnleanIssueBase#1}{GitHub issue \##1}}
\newcommand{\ghpr}[1]{\href{\tnleanPrBase#1}{PR \##1}}
\newcommand{\doclink}[1]{\href{\tnleanDocBase#1.html}{\path{#1}}}

% Dirac notation and the identity operator, as in blueprint/src/macros/common.tex.
% \providecommand keeps note-local definitions authoritative.
\usepackage{dsfont}
\providecommand{\ket}[1]{|#1\rangle}
\providecommand{\bra}[1]{\langle #1|}
\providecommand{\braket}[2]{\langle #1 | #2 \rangle}
\providecommand{\ketbra}[2]{|#1\rangle\!\langle #2|}
\providecommand{\Id}{\mathds{1}}
\providecommand{\one}{\mathds{1}}
\providecommand{\C}{\mathbb{C}}
\providecommand{\MN}[1]{M_{#1}(\C)}
\providecommand{\MD}{M_D(\C)}

% Verdict marker: \gapnote{<kind>}{<status>}. Kind is the policy
% classification (clarification, local-correction, scope-restriction,
% unfaithful, false-source, open-gap); status is open, wip (elimination
% underway), resolved, or historical. Machine-read by the site index;
% prints a verdict line.
\newcommand{\gapnote}[2]{\par\noindent\textbf{Verdict:} #1 (#2).\par}

\providecommand{\gapnote}[2]{\par\noindent\textbf{Verdict:} #1 (#2).\par}
\providecommand{\leanid}[1]{\textsf{\detokenize{#1}}}
\title{Bond Dimension Zero in the Viable Sets of Landau, Vazirani and Vidick}
\author{}
\date{}
\begin{document}
\maketitle
\gapnote{local-correction}{resolved}

\section{The source statement}
Landau, Vazirani and Vidick~\cite{Landau2015Polynomial},
\href{https://arxiv.org/abs/1307.5143v1}{arXiv:1307.5143v1}, source
\path{algo53.tex}, line 243, call a set $S$ of vectors on the first $i$ sites
$(i,s,b,\delta)$-viable if, among other conditions, every vector of $S$ is a
matrix product state of bond dimension at most $b$. In the proof of Claim
\texttt{c:ext} (line 399) they observe that appending one site in a basis state
does not raise the bond dimension.

\section{The point at issue}
Read literally, the observation fails at $b=0$ under the library's notion of an
open matrix product state, in which a vector on one site has no internal bond
and so has bond dimension $0$, while a nonzero vector on two sites has an
internal bond of dimension at least $1$. Appending a basis state to the vector
$\lvert0\rangle$ on one site gives the nonzero product vector
$\lvert0\rangle\lvert k\rangle$ on two sites, which has bond dimension $1$, not
$0$. For $b\ge1$ the observation holds: the new internal bond carries a vector
of rank one.

\section{The convention}
The bond dimensions in the source are positive polynomials in the number of
sites, so $b=0$ never occurs. The formalization reads ``bond dimension at most
$b$'' as ``bond dimension at most $\max(1,b)$'', in the definition of a matrix
product state on the first sites, so no theorem carries a hypothesis $b>0$.

\section{Formalization}
The definition is \leanid{TNLean.MPS.AreaLaw.IsLeftOpenMPS}, and the extension
step is \leanid{TNLean.MPS.AreaLaw.IsLeftOpenMPS.extendSite} and
\leanid{TNLean.MPS.AreaLaw.IsViableMPS.extend}.
\bibliographystyle{unsrt}
\bibliography{references}
\end{document}
